\documentclass[a4paper]{article}

% Basic packages
% Español (México) con XeLaTeX: fontspec para fuentes Unicode, polyglossia
% para la división silábica y los nombres del documento en la variante mexicana.
\usepackage{fontspec}
\usepackage{polyglossia}
\setdefaultlanguage[variant=mexican]{spanish}
% Los nombres chinos van como «pinyin (original)»: xeCJK da a los caracteres
% Han su propia fuente; sin ella XeLaTeX los omite con un aviso.
% FandolSong no cubre algunos caracteres de nombres (U+73FA); WenQuanYi Zen Hei sí.
\usepackage[AutoFallBack=true]{xeCJK}
\setCJKmainfont{FandolSong-Regular.otf}
\setCJKfallbackfamilyfont{\CJKrmdefault}{WenQuanYi Zen Hei}
% polyglossia no traduce los nombres de listings.
\AtBeginDocument{\providecommand{\lstlistingname}{}\renewcommand{\lstlistingname}{Listado}%
  \providecommand{\lstlistlistingname}{}\renewcommand{\lstlistlistingname}{Índice de listados}}
\usepackage{amsmath, amssymb}
\usepackage{graphicx}
\usepackage[margin=2.5cm]{geometry}
\usepackage[most]{tcolorbox}
\usepackage{etoolbox}
\usepackage{listings}
\usepackage{hyperref}
\usepackage{booktabs}
\usepackage{subcaption}
\usepackage{float}
\usepackage{tikz}
\IfFileExists{pgfplots.sty}{
    \usepackage{pgfplots}
    \pgfplotsset{compat=1.18}
}{}

% Highlight boxes
\newtcolorbox{knowledgebox}[1]{
    enhanced, colback=blue!5!white, colframe=blue!75!black, colbacktitle=blue!75!black,
    coltitle=white, fonttitle=\bfseries, title={#1}, attach boxed title to top left={yshift=-2mm, xshift=2mm},
    boxrule=1pt, sharp corners
}

\newtcolorbox{importantbox}[1]{
    enhanced, colback=yellow!10!white, colframe=yellow!80!black, colbacktitle=yellow!80!black,
    coltitle=black, fonttitle=\bfseries, title={#1}, sharp corners
}

\newtcolorbox{warningbox}[1]{
    enhanced, colback=red!5!white, colframe=red!75!black, colbacktitle=red!75!black,
    coltitle=white, fonttitle=\bfseries, title={#1}, sharp corners
}

% Listings
\lstset{
    language=Python,
    basicstyle=\ttfamily\small,
    keywordstyle=\color{blue},
    stringstyle=\color{red!60!black},
    commentstyle=\color{green!60!black},
    breaklines=true,
    frame=single,
    numbers=left,
    numberstyle=\tiny\color{gray},
    captionpos=b,
    extendedchars=true
}

% Front-page metadata
\newcommand{\notetitle}{MIT 6.S191 Lecture 2: Recurrent Neural Networks, Transformers, and Attention}
\newcommand{\noteauthors}{Elaboradas a partir de la clase de Ava Amini}
\newcommand{\notedate}{Primavera de 2025}
\newcommand{\videochannel}{Alexander Amini (MIT)}
\newcommand{\videopublishdate}{2025-03-10}
\newcommand{\videoduration}{61 minutos}
\newcommand{\videourl}{https://www.youtube.com/watch?v=GvezxUdLrEk}
\newcommand{\videocoverpath}{cover.jpg}
\newcommand{\repourl}{https://github.com/hqhq1025/ai-course-notes}

% Footer with repo info
\usepackage{fancyhdr}
\pagestyle{fancy}
\fancyhf{}
\fancyfoot[C]{\thepage}
\fancyfoot[R]{\footnotesize\color{gray}\href{\repourl}{AI Course Notes}}
\renewcommand{\headrulewidth}{0pt}
\renewcommand{\footrulewidth}{0pt}

\begin{document}

\begin{titlepage}
\centering
{\Large Notas del curso\par}
\vspace{1.2cm}
{\huge\bfseries \notetitle\par}
\vspace{0.8cm}
{\large \noteauthors\par}
\vspace{0.3cm}
{\large \notedate\par}
\vspace{1.2cm}

\ifdefempty{\videocoverpath}{
{\small Al generar las notas, indique la ruta local de la portada del video.\par}
}{
\includegraphics[width=0.82\textwidth,height=0.45\textheight,keepaspectratio]{\videocoverpath}\par
}

\vfill
\begin{tcolorbox}[width=0.9\textwidth, colback=black!2!white, colframe=black!60, sharp corners]
\textbf{Autor o canal del video}: \videochannel\par
\textbf{Fecha de publicación}: \videopublishdate\par
\textbf{Duración del video}: \videoduration\par
\ifdefempty{\videourl}{}{
\textbf{Enlace del video}: \href{\videourl}{\nolinkurl{\videourl}}\par
}\par
\textbf{Página del proyecto}: \href{\repourl}{\nolinkurl{\repourl}}\par
\end{tcolorbox}
\end{titlepage}

\tableofcontents
\newpage

%% --- Inicio del contenido --- %%

\section{Introducción: por qué se necesita el modelado de secuencias}

Esta clase es la segunda de MIT 6.S191 (Introduction to Deep Learning), impartida por Ava Amini. En la clase anterior, Alexander Amini presentó los fundamentos de las redes neuronales, incluyendo el perceptrón (Perceptron), la red de propagación hacia adelante (Feedforward Network) y el algoritmo de retropropagación (Backpropagation). A partir de estos fundamentos, esta clase explica de manera sistemática cómo aplicar las redes neuronales al procesamiento de \textbf{datos secuenciales} (Sequential Data), y abarca tres temas centrales:

\begin{enumerate}
    \item \textbf{Recurrent Neural Networks (RNNs)}: los principios básicos y el mecanismo interno de las redes neuronales recurrentes
    \item \textbf{Los retos de entrenar una RNN}: el problema del desvanecimiento/explosión del gradiente y arquitecturas mejoradas como la LSTM
    \item \textbf{El mecanismo de Attention y el Transformer}: el principio del mecanismo de autoatención y su aplicación en los modelos de lenguaje grandes modernos
\end{enumerate}

\begin{importantbox}{La idea central del modelado de secuencias}
El objetivo central del modelado de secuencias es: dada la información histórica de una secuencia, predecir su estado futuro. Esta es la base de numerosas aplicaciones de IA, como los modelos de lenguaje grandes (LLM), el reconocimiento de voz y la traducción automática.
\end{importantbox}

\subsection{Introducción intuitiva: predecir el movimiento de una pelota}

La ponente introduce el concepto de modelado de secuencias con un ejemplo ilustrativo: supongamos que hay una pelota en un espacio bidimensional, y nuestra tarea es predecir la dirección de su próximo movimiento.

\begin{figure}[H]
\centering
\includegraphics[width=0.7\textwidth]{slides-images/slide-03.jpg}
\caption{Solo con la posición actual de la pelota no se puede determinar la dirección de su movimiento\protect\footnotemark}
\end{figure}
\footnotetext{Diapositivas, página 3.}

Si solo se da la posición actual de la pelota, sin ninguna información histórica, nuestra predicción solo puede ser aleatoria: la pelota podría moverse en cualquier dirección. Pero si se da la secuencia de posiciones de la pelota en los pasos de tiempo anteriores, el problema se vuelve mucho más sencillo:

\begin{figure}[H]
\centering
\includegraphics[width=0.7\textwidth]{slides-images/slide-05.jpg}
\caption{Usar la secuencia de posiciones históricas permite predecir con precisión la dirección del movimiento\protect\footnotemark}
\end{figure}
\footnotetext{Diapositivas, página 5.}

Al observar la trayectoria histórica de la pelota de izquierda a derecha, podemos inferir razonablemente que seguirá moviéndose hacia la derecha. Este ejemplo revela la idea central del modelado de secuencias: \textbf{usar la información pasada para predecir el futuro}.

\subsection{Los datos secuenciales están en todas partes}

Los datos secuenciales están presentes en todas partes del mundo real:

\begin{figure}[H]
\centering
\includegraphics[width=0.85\textwidth]{slides-images/slide-07.jpg}
\caption{Tipos de datos secuenciales en el mundo real\protect\footnotemark}
\end{figure}
\footnotetext{Diapositivas, página 7.}

\begin{itemize}
    \item \textbf{Señales de audio}: la forma de onda del habla puede dividirse en una serie de segmentos de onda sonora para procesarse como secuencia
    \item \textbf{Lenguaje natural}: el texto puede dividirse en una secuencia de palabras o caracteres
    \item \textbf{Señales médicas}: por ejemplo, la serie de tiempo de un electrocardiograma (ECG)
    \item \textbf{Datos financieros}: la serie de tiempo del precio de las acciones
    \item \textbf{Secuencias biológicas}: secuencias de ácidos nucleicos de ADN/ARN, secuencias de aminoácidos de proteínas
    \item \textbf{Datos meteorológicos, secuencias de fotogramas de video}, entre otros
\end{itemize}

\subsection{Tipos de tareas del modelado de secuencias}

\begin{knowledgebox}{Los tres patrones básicos del modelado de secuencias}
\begin{itemize}
    \item \textbf{Many-to-One}: entrada de secuencia $\rightarrow$ salida única. Por ejemplo: la clasificación de sentimiento (se ingresa un fragmento de texto y se produce positivo/negativo)
    \item \textbf{One-to-Many}: entrada única $\rightarrow$ salida de secuencia. Por ejemplo: la generación de descripciones de imágenes (se ingresa una imagen y se produce el texto descriptivo)
    \item \textbf{Many-to-Many}: entrada de secuencia $\rightarrow$ salida de secuencia. Por ejemplo: la traducción automática (secuencia en inglés $\rightarrow$ secuencia en chino)
\end{itemize}
\end{knowledgebox}

\begin{figure}[H]
\centering
\includegraphics[width=0.75\textwidth]{slides-images/slide-08.jpg}
\caption{Los distintos tipos de tareas del modelado de secuencias\protect\footnotemark}
\end{figure}
\footnotetext{Diapositivas, página 8.}

\subsection{Resumen de la sección}

El modelado de secuencias es uno de los paradigmas más importantes del aprendizaje profundo, y su núcleo consiste en aprovechar la dependencia temporal/secuencial de los datos. Desde la simple predicción del movimiento de una pelota hasta los modelos de lenguaje complejos, la idea del modelado de secuencias atraviesa todos los aspectos de la IA moderna.

\newpage
\section{Del perceptrón a las redes neuronales recurrentes}

\subsection{Repaso del perceptrón}

En la clase anterior, aprendimos la estructura básica del perceptrón: dados los insumos $x^{(1)}, x^{(2)}, \ldots, x^{(m)}$, cada insumo se multiplica por el peso correspondiente $w_1, w_2, \ldots, w_m$, se suman y el resultado pasa por una función de activación no lineal para producir la salida.

\begin{figure}[H]
\centering
\includegraphics[width=0.7\textwidth]{slides-images/slide-10.jpg}
\caption{Repaso de la estructura del perceptrón\protect\footnotemark}
\end{figure}
\footnotetext{Página 10 de las diapositivas.}

Expresado con fórmulas matemáticas:

$$z = \sum_{i=1}^{m} w_i x^{(i)} + b$$

$$\hat{y} = g(z)$$

Donde:
\begin{itemize}
    \item $x^{(i)}$: la $i$-ésima característica de entrada
    \item $w_i$: el peso correspondiente
    \item $b$: el término de sesgo
    \item $g(\cdot)$: la función de activación no lineal
    \item $\hat{y}$: la salida predicha
\end{itemize}

\subsection{Limitaciones del perceptrón para procesar datos secuenciales}

Cuando aplicamos el perceptrón (o una red feedforward) directamente a datos secuenciales, nos encontramos con un problema fundamental: \textbf{la predicción de cada paso de tiempo es completamente independiente}, sin aprovechar la información de los demás pasos de tiempo.

\begin{figure}[H]
\centering
\includegraphics[width=0.75\textwidth]{slides-images/slide-13.jpg}
\caption{Una red feedforward procesa cada paso de tiempo de forma independiente: no puede aprovechar la información histórica\protect\footnotemark}
\end{figure}
\footnotetext{Página 13 de las diapositivas.}

En concreto, si usamos una red feedforward para procesar cada paso de tiempo de forma independiente:

$$\hat{y}_t = f(x_t)$$

Entonces la predicción $\hat{y}_2$ del paso de tiempo $t=2$ depende únicamente del insumo $x_2$, ignorando por completo la información clave que podrían contener $x_0$ y $x_1$.

Esto es como observar solo la posición actual de una pelota sin tomar en cuenta su trayectoria anterior: la predicción es necesariamente ciega.

\subsection{La introducción de la recurrencia: la idea central del Recurrence}

Para resolver este problema, necesitamos un mecanismo que transmita al paso de tiempo actual la información de los pasos de tiempo pasados. La idea central es:

\begin{importantbox}{El núcleo de las redes neuronales recurrentes: el estado oculto}
Se define un \textbf{estado interno (Hidden State)} $h_t$, que se transmite entre los pasos de tiempo y conserva la ``memoria'' de la red sobre los insumos pasados. La predicción de cada paso de tiempo depende no solo del insumo actual $x_t$, sino también del estado oculto del paso anterior $h_{t-1}$:

$$\hat{y}_t = f(x_t, h_{t-1})$$
\end{importantbox}

\begin{figure}[H]
\centering
\includegraphics[width=0.85\textwidth]{slides-images/slide-15.jpg}
\caption{Una neurona recurrente: la salida depende del insumo actual y de la memoria pasada\protect\footnotemark}
\end{figure}
\footnotetext{Página 15 de las diapositivas.}

La imagen anterior muestra a la izquierda la vista plegada (Folded View) de la RNN y a la derecha la vista desplegada (Unrolled View). El estado oculto $h_t$ se transmite entre los pasos de tiempo como si fuera una cadena, lo que permite a la red ``recordar'' la información que vio previamente.

\subsection{Resumen de la sección}

La trayectoria del perceptrón a la RNN es clara: las redes feedforward carecen de la capacidad de procesar secuencias, por lo que introducimos el estado oculto como puente de información entre los pasos de tiempo, dando forma al concepto de recurrencia (Recurrence). Esta es la base para entender todos los modelos de secuencias.

\newpage
\section{Mecanismo interno de las RNN}

\subsection{Intuición de las RNN}

El docente construye la intuición de las RNN mediante un fragmento de pseudocódigo en Python:

\begin{figure}[H]
\centering
\includegraphics[width=0.85\textwidth]{slides-images/slide-20.jpg}
\caption{Intuición del pseudocódigo de una RNN: se recorre cada palabra de la oración, se actualiza el estado oculto y se genera una predicción\protect\footnotemark}
\end{figure}
\footnotetext{Diapositivas, página 20.}

\begin{lstlisting}[caption={Ejemplo de pseudocódigo de una RNN}]
my_rnn = RNN()
hidden_state = [0, 0, 0, 0]

sentence = ["I", "love", "recurrent", "neural"]

for word in sentence:
    prediction, hidden_state = my_rnn(word, hidden_state)

next_word_prediction = prediction
# >>> "networks!"
\end{lstlisting}

Este pseudocódigo muestra con claridad el flujo de trabajo de una RNN:
\begin{enumerate}
    \item Se inicializa el estado oculto como un vector de ceros
    \item Para cada elemento de la secuencia, se invoca la RNN con la entrada actual y el estado oculto del paso anterior
    \item Cada invocación devuelve una predicción y el estado oculto actualizado
    \item Una vez procesada toda la secuencia, la predicción final es la salida
\end{enumerate}

\subsection{Fórmula matemática de la actualización de estado de una RNN}

\begin{figure}[H]
\centering
\includegraphics[width=0.75\textwidth]{slides-images/slide-22.jpg}
\caption{Mecanismo de actualización de estado y salida de una RNN\protect\footnotemark}
\end{figure}
\footnotetext{Diapositivas, página 22.}

La propagación hacia adelante de una RNN comprende dos pasos clave:

\textbf{Paso 1: actualizar el estado oculto}

$$h_t = \tanh(W_{xh} \cdot x_t + W_{hh} \cdot h_{t-1} + b_h)$$

\begin{itemize}
    \item $W_{xh}$: matriz de pesos que transforma la entrada al espacio oculto
    \item $W_{hh}$: matriz de pesos de la actualización del propio estado oculto
    \item $b_h$: término de sesgo
    \item $\tanh$: función de activación no lineal
\end{itemize}

\textbf{Paso 2: generar la salida}

$$\hat{y}_t = W_{hy} \cdot h_t + b_y$$

\begin{itemize}
    \item $W_{hy}$: matriz de pesos que transforma el estado oculto en la salida
    \item $b_y$: sesgo de la salida
\end{itemize}

\begin{importantbox}{Compartición de pesos: la característica clave de las RNN}
En una RNN, las matrices de pesos $W_{xh}$, $W_{hh}$ y $W_{hy}$ se comparten en \textbf{todos los pasos de tiempo}. Esto significa que, sin importar la longitud de la secuencia, el número de parámetros no cambia. Esto no solo reduce en gran medida los parámetros del modelo, sino que además permite que la RNN procese secuencias de cualquier longitud.
\end{importantbox}

\subsection{Despliegue del grafo de cómputo de una RNN}

\begin{figure}[H]
\centering
\includegraphics[width=0.85\textwidth]{slides-images/slide-25.jpg}
\caption{Despliegue del grafo de cómputo de una RNN a lo largo del eje temporal\protect\footnotemark}
\end{figure}
\footnotetext{Diapositivas, página 25.}

Al desplegar la RNN a lo largo del eje temporal, se aprecia con claridad:
\begin{itemize}
    \item Cada paso de tiempo recibe la entrada $x_t$, que participa en el cálculo del estado oculto a través de $W_{xh}$
    \item El estado oculto $h_t$ se transmite al siguiente paso de tiempo a través de $W_{hh}$
    \item Cada paso de tiempo produce la predicción de salida $\hat{y}_t$ a través de $W_{hy}$
\end{itemize}

\subsection{Implementación en código de una RNN}

\begin{lstlisting}[caption={Implementación de una capa RNN con TensorFlow}]
class MyRNNCell(tf.keras.layers.Layer):
    def __init__(self, rnn_units):
        super(MyRNNCell, self).__init__()
        self.W_xh = self.add_weight([rnn_units, rnn_units])
        self.W_hh = self.add_weight([rnn_units, rnn_units])
        self.W_hy = self.add_weight([rnn_units, num_outputs])

    def call(self, x, h):
        # Update hidden state
        h = tf.math.tanh(
            tf.linalg.matmul(x, self.W_xh) +
            tf.linalg.matmul(h, self.W_hh)
        )
        # Compute output
        output = tf.linalg.matmul(h, self.W_hy)
        return output, h
\end{lstlisting}

En marcos de trabajo como TensorFlow/Keras, también se puede usar directamente una capa RNN ya encapsulada:

\begin{lstlisting}[caption={Uso de la capa SimpleRNN encapsulada de Keras}]
tf.keras.layers.SimpleRNN(rnn_units)
\end{lstlisting}

\subsection{Resumen de la sección}

La operación central de una RNN puede resumirse en tres pasos: (1) recibir la entrada actual; (2) combinarla con el estado oculto del paso anterior para actualizar el estado oculto actual; (3) producir una salida a partir del estado oculto. La compartición de pesos en todos los pasos de tiempo es la clave que permite a la RNN procesar secuencias de longitud variable.

\newpage
\section{Aplicación práctica del modelado de secuencias: modelado del lenguaje}

\subsection{La tarea de Next Word Prediction}

El modelado del lenguaje (Language Modeling) es una de las aplicaciones más representativas e importantes del modelado de secuencias. Su tarea central es extremadamente simple: \textbf{dada una secuencia de palabras, predecir la siguiente palabra más probable}.

\begin{figure}[H]
\centering
\includegraphics[width=0.85\textwidth]{slides-images/slide-35.jpg}
\caption{Predecir la siguiente palabra: la tarea central del modelado del lenguaje\protect\footnotemark}
\end{figure}
\footnotetext{Diapositiva 35.}

Por ejemplo, dada la oración ``This morning I took my cat for a'', el modelo necesita predecir que la siguiente palabra es ``walk''.

\begin{importantbox}{Next Word Prediction es la base de los LLM}
Los grandes modelos de lenguaje actuales (como la serie GPT) están, en esencia, ejecutando esta tarea de Next Word Prediction. Al entrenarse sobre enormes cantidades de texto para predecir la siguiente palabra, el modelo aprende la estructura y la semántica del lenguaje.
\end{importantbox}

\subsection{Cómo representar el lenguaje mediante números}

Las redes neuronales solo pueden procesar datos numéricos, por lo que necesitamos convertir las palabras en vectores numéricos. Este proceso consta de dos pasos:

\textbf{Paso 1: Construir el vocabulario (Vocabulary)}

Definir un vocabulario finito que contenga todas las palabras posibles y asignar a cada palabra un número de índice único:
\begin{itemize}
    \item ``a'' $\rightarrow$ 1
    \item ``cat'' $\rightarrow$ 2
    \item ``the'' $\rightarrow$ 3
    \item $\ldots$
\end{itemize}

\textbf{Paso 2: Embedding}

Mapear los índices a vectores de dimensión fija. El método más simple es \textbf{One-Hot Encoding}:

$$\text{``cat''} \rightarrow [0, 1, 0, 0, \ldots, 0]$$

\begin{warningbox}{Las limitaciones del One-Hot Encoding}
El One-Hot Encoding tiene dos defectos graves:
\begin{enumerate}
    \item \textbf{La maldición de la dimensionalidad}: la dimensión del vector es igual al tamaño del vocabulario; para un vocabulario con decenas o incluso cientos de miles de palabras, esta representación resulta extremadamente dispersa e ineficiente
    \item \textbf{Ausencia de información semántica}: la ``distancia'' entre todas las palabras es igual; la distancia entre ``cat'' y ``dog'' es la misma que entre ``cat'' y ``airplane'', por lo que no puede capturar la similitud semántica
\end{enumerate}
\end{warningbox}

Un método mejor es usar \textbf{Learned Embedding} (embedding aprendido): mediante entrenamiento, la red aprende de manera automática una representación vectorial de baja dimensión y densa, de modo que las palabras semánticamente cercanas queden a menor distancia en el espacio vectorial.

\begin{figure}[H]
\centering
\includegraphics[width=0.75\textwidth]{slides-images/slide-38.jpg}
\caption{Learned Embedding: las palabras semánticamente cercanas quedan a menor distancia en el espacio vectorial\protect\footnotemark}
\end{figure}
\footnotetext{Diapositiva 38.}

\begin{knowledgebox}{La naturaleza del embedding}
Un embedding es, en esencia, una tabla de consulta (Lookup Table) que puede aprenderse. Mapea símbolos discretos (como palabras o caracteres) a un espacio vectorial continuo. Después del entrenamiento, las palabras semánticamente similares (como ``king'' y ``queen'') quedan cercanas entre sí en el espacio de embedding, mientras que las palabras con significados distintos quedan más alejadas. Esta propiedad resulta fundamental para las tareas posteriores de modelado de secuencias.
\end{knowledgebox}

\subsection{Dependencias de corto y largo plazo en las secuencias}

En el lenguaje natural existen relaciones de dependencia de distinto alcance:

\begin{itemize}
    \item \textbf{Dependencia de corto plazo}: ``The clouds are in the \underline{\hspace{1cm}}'' $\rightarrow$ ``sky''. La respuesta depende solo de las palabras inmediatamente anteriores.
    \item \textbf{Dependencia de largo plazo}: ``I grew up in France, $\ldots$ and I speak fluent \underline{\hspace{1cm}}'' $\rightarrow$ ``French''. La respuesta depende de un contexto muy anterior, ``France''.
\end{itemize}

Las dependencias de largo plazo plantean un desafío severo para las RNN, lo cual da paso directamente al tema de la siguiente sección.

\subsection{Resumen de la sección}

La tarea central del modelado del lenguaje —predecir la siguiente palabra— parece simple pero es extremadamente poderosa: es la base de entrenamiento de todos los grandes modelos de lenguaje modernos. Representar el lenguaje mediante vectores numéricos (embedding) es el primer paso para lograr esta tarea, mientras que manejar las dependencias de corto y largo plazo en la secuencia constituye el desafío central.

\newpage
\section{Entrenamiento de RNN: retropropagación y problemas de gradiente}

\subsection{Función de pérdida de la RNN}

Al igual que en las redes feedforward, entrenar una RNN también requiere definir una función de pérdida. La diferencia es que la RNN puede calcular un valor de pérdida en \textbf{cada paso de tiempo}:

$$\mathcal{L}_t = \text{Loss}(\hat{y}_t, y_t)$$

La pérdida total es la suma de las pérdidas de todos los pasos de tiempo:

$$\mathcal{L} = \sum_{t=1}^{T} \mathcal{L}_t$$

\subsection{Backpropagation Through Time (BPTT)}

Entrenar una RNN utiliza una versión extendida de la retropropagación: la \textbf{retropropagación a través del tiempo} (Backpropagation Through Time, BPTT). Su idea central es:

\begin{enumerate}
    \item Propagación hacia adelante: calcular el estado oculto y la salida paso a paso en el tiempo
    \item Calcular la pérdida en cada paso de tiempo
    \item Retropropagación: propagar el error a lo largo del eje temporal de atrás hacia adelante, calculando el gradiente de cada parámetro
    \item Actualizar los parámetros
\end{enumerate}

\begin{figure}[H]
\centering
\includegraphics[width=0.85\textwidth]{slides-images/slide-45.jpg}
\caption{Flujo del gradiente en una RNN estándar\protect\footnotemark}
\end{figure}
\footnotetext{Diapositiva 45.}

\subsection{Desvanecimiento y explosión del gradiente}

Durante el proceso de BPTT, el gradiente debe pasar por múltiples multiplicaciones de matrices (multiplicaciones sucesivas de $W_{hh}$) y por el producto sucesivo de las derivadas de la función de activación. Esto provoca dos problemas numéricos graves:

\begin{importantbox}{Desvanecimiento y explosión del gradiente}
\begin{itemize}
    \item \textbf{Explosión del gradiente} (Exploding Gradients): cuando el valor propio de $W_{hh}$ es $> 1$, el gradiente crece exponencialmente con los pasos de tiempo, lo que produce actualizaciones de parámetros excesivas e inestabilidad en el entrenamiento
    \item \textbf{Desvanecimiento del gradiente} (Vanishing Gradients): cuando el valor propio de $W_{hh}$ es $< 1$, el gradiente decae exponencialmente con los pasos de tiempo hasta acercarse a cero, lo que impide que la red aprenda dependencias de largo plazo
\end{itemize}
\end{importantbox}

\begin{figure}[H]
\centering
\includegraphics[width=0.85\textwidth]{slides-images/slide-48.jpg}
\caption{El problema del desvanecimiento del gradiente: calcular el gradiente respecto de $h_0$ implica multiplicaciones sucesivas de $W_{hh}$\protect\footnotemark}
\end{figure}
\footnotetext{Diapositiva 48.}

\subsection{La dificultad de las dependencias de largo plazo}

El problema del desvanecimiento del gradiente hace que a la RNN le resulte difícil capturar dependencias de largo plazo:

\begin{figure}[H]
\centering
\includegraphics[width=0.85\textwidth]{slides-images/slide-52.jpg}
\caption{El desvanecimiento del gradiente hace que la RNN tienda a capturar dependencias de corto plazo\protect\footnotemark}
\end{figure}
\footnotetext{Diapositiva 52.}

\begin{warningbox}{Desvanecimiento del gradiente $\neq$ gradiente igual a cero}
El desvanecimiento del gradiente no significa que el gradiente desaparezca por completo, sino que la señal de gradiente proveniente de \textbf{pasos de tiempo más lejanos} se vuelve extremadamente débil, de modo que la actualización de los parámetros del modelo queda dominada casi por completo por el gradiente de los pasos de tiempo recientes. Por ello, el modelo tiende a aprender dependencias de corto plazo y le cuesta capturar patrones de largo plazo.
\end{warningbox}

\subsection{Estrategias para resolver los problemas de gradiente}

Existen tres grandes tipos de métodos para resolver el desvanecimiento/explosión del gradiente:

\begin{enumerate}
    \item \textbf{Elección de la función de activación}: usar funciones de activación como ReLU, que no comprimen el gradiente en la zona de saturación
    \item \textbf{Inicialización de los pesos}: diseñar cuidadosamente la estrategia de inicialización de los pesos para evitar que el gradiente se escale rápidamente durante la propagación
    \item \textbf{Mejoras en la arquitectura de red}: diseñar mecanismos de compuertas específicos para controlar el flujo de información, como LSTM y GRU
\end{enumerate}

\subsection{LSTM: redes de memoria de largo y corto plazo}

\textbf{Long Short-Term Memory (LSTM)} es la variante de RNN más conocida para resolver el problema del desvanecimiento del gradiente. Su idea central es agregar \textbf{mecanismos de compuertas} (Gates) dentro de la unidad estándar de la RNN, con los que se controla de manera selectiva el olvido, el almacenamiento y la salida de información.

\begin{importantbox}{Las tres compuertas de la LSTM}
\begin{itemize}
    \item \textbf{Compuerta de olvido} (Forget Gate): decide qué información antigua debe descartarse
    \item \textbf{Compuerta de entrada} (Input Gate): decide qué información nueva debe almacenarse
    \item \textbf{Compuerta de salida} (Output Gate): decide qué información debe emitirse hacia el siguiente paso de tiempo
\end{itemize}
El valor de estas compuertas se controla mediante la función sigmoide, cuya salida está entre $[0, 1]$ y representa desde «olvido completo» hasta «retención completa».
\end{importantbox}

Mediante este mecanismo de compuertas, la LSTM permite que el gradiente se propague entre pasos de tiempo a lo largo de una especie de «autopista» (Cell State), lo que alivia en gran medida el problema del desvanecimiento del gradiente y permite que la red aprenda dependencias de más largo plazo.

\begin{knowledgebox}{Importancia histórica de la LSTM}
La LSTM fue propuesta por Hochreiter y Schmidhuber en 1997, y es una de las innovaciones arquitectónicas más importantes en la historia del aprendizaje profundo. Antes de la aparición del Transformer (en 2017), la LSTM y sus variantes (como la GRU) fueron la opción dominante para procesar datos secuenciales, con un uso extendido en la traducción automática, el reconocimiento de voz, la generación de texto y otras áreas.
\end{knowledgebox}

\subsection{Resumen de la sección}

El reto central al entrenar una RNN es el problema del desvanecimiento/explosión del gradiente, que limita la capacidad de la RNN estándar para aprender dependencias de largo plazo. La LSTM alivia eficazmente este problema mediante su mecanismo de compuertas, pero en esencia sigue procesando la secuencia paso a paso en el tiempo, con limitaciones tanto en la eficiencia computacional como en la capacidad de modelar secuencias largas.

\newpage
\section{Aplicaciones y limitaciones de las RNN}

\subsection{Un caso de aplicación de las RNN: generación de música}

El presentador presentó la aplicación de las RNN en la generación de música. Esta es una tarea naturalmente adecuada para el modelado de secuencias: dada la secuencia histórica de notas musicales, se predice la nota siguiente más probable.

\begin{figure}[H]
\centering
\includegraphics[width=0.85\textwidth]{slides-images/slide-56.jpg}
\caption{Generación de música con RNN: predicción de la siguiente nota\protect\footnotemark}
\end{figure}
\footnotetext{Diapositivas, página 56.}

El presentador mencionó un caso interesante: un equipo entrenó un modelo de red neuronal para aprender música clásica y después le pidió al modelo que intentara completar el tercer movimiento de la célebre «Sinfonía inconclusa» de Franz Schubert. Esto muestra el potencial de las RNN en tareas creativas.

\subsection{Criterios de diseño del modelado de secuencias}

El presentador resumió cuatro criterios de diseño clave del modelado de secuencias:

\begin{figure}[H]
\centering
\includegraphics[width=0.6\textwidth]{slides-images/slide-30.jpg}
\caption{Criterios de diseño del modelado de secuencias\protect\footnotemark}
\end{figure}
\footnotetext{Diapositivas, página 30.}

\begin{enumerate}
    \item \textbf{Manejar secuencias de longitud variable} (Variable-length sequences)
    \item \textbf{Rastrear dependencias de largo plazo} (Long-term dependencies)
    \item \textbf{Mantener información sobre el orden} (Maintain information about order)
    \item \textbf{Compartir parámetros a lo largo de la secuencia} (Share parameters across the sequence)
\end{enumerate}

Las RNN satisfacen estos cuatro criterios, pero en la práctica siguen teniendo limitaciones importantes.

\subsection{Limitaciones centrales de las RNN}

\begin{warningbox}{Las tres grandes limitaciones de las RNN}
\begin{enumerate}
    \item \textbf{Cuello de botella del cómputo secuencial}: las RNN deben procesar los datos paso a paso en el tiempo y no pueden paralelizarse. Para secuencias largas, tanto el entrenamiento como la inferencia son muy lentos
    \item \textbf{Dificultad con las dependencias de largo plazo}: aunque las LSTM mitigan el problema del desvanecimiento del gradiente, para secuencias muy largas la información sigue atenuándose de manera gradual
    \item \textbf{Cuello de botella de codificación}: en la arquitectura Encoder-Decoder, la información de toda la secuencia de entrada debe comprimirse en un vector de estado oculto de tamaño fijo, lo cual resulta insuficiente para secuencias largas
\end{enumerate}
\end{warningbox}

Estas limitaciones llevaron a los investigadores a preguntarse: \textbf{¿es posible prescindir por completo del mecanismo de recurrencia (Recurrence) y procesar los datos de secuencia de una manera totalmente nueva?}

\subsection{Resumen de la sección}

Las RNN (incluidas las LSTM/GRU) lograron un gran éxito en el modelado de secuencias, pero su naturaleza de procesamiento secuencial limita la eficiencia computacional, y las dificultades en la propagación del gradiente también restringen su capacidad para manejar secuencias extremadamente largas. Estas limitaciones impulsaron directamente el surgimiento del mecanismo de Attention y de la arquitectura Transformer.

\newpage
\section{Mecanismo de Attention: la atención es todo lo que necesitas}

\subsection{Más allá de la recurrencia: en busca de un nuevo paradigma}

\begin{figure}[H]
\centering
\includegraphics[width=0.85\textwidth]{slides-images/slide-60.jpg}
\caption{El objetivo del modelado de secuencias: de la secuencia de entrada a las características y de ahí a la salida\protect\footnotemark}
\end{figure}
\footnotetext{Diapositiva 60.}

Volvamos al objetivo fundamental del modelado de secuencias: dada una secuencia de entrada, extraer características y generar una predicción. Las RNN capturan las dependencias de la secuencia procesándola paso a paso, pero ¿podríamos hacerlo mejor si pudiéramos \textbf{ver toda la secuencia a la vez} e identificar automáticamente las partes importantes?

Una idea ingenua sería concatenar todas las entradas en un solo vector grande y procesarlo con una red feedforward. Pero esto presenta dos problemas:
\begin{itemize}
    \item El número de parámetros crece de manera explosiva con la longitud de la secuencia, lo cual no es escalable
    \item La operación de concatenación destruye la información de posición y de orden
\end{itemize}

\begin{figure}[H]
\centering
\includegraphics[width=0.75\textwidth]{slides-images/slide-62.jpg}
\caption{El problema del método ingenuo: concatenar las entradas pierde la información de posición\protect\footnotemark}
\end{figure}
\footnotetext{Diapositiva 62.}

\subsection{Qué es Attention}

La idea central del mecanismo de \textbf{Attention} (atención) proviene de la cognición humana: cuando observamos una escena, no prestamos atención de manera uniforme a todo su contenido, sino que nos enfocamos automáticamente en las partes más relevantes e importantes.

\begin{importantbox}{Definición central de Attention}
Attention es un mecanismo que permite a la red \textbf{identificar automáticamente y enfocarse en las partes más importantes de la entrada}. Calcula \textbf{puntuaciones de relevancia} (Attention Score) entre las distintas partes de la entrada, para decidir a cuáles de ellas debe «prestar atención» al generar la salida.
\end{importantbox}

\subsection{Analogía con un motor de búsqueda: Query-Key-Value}

Quien dio la clase usa una analogía muy intuitiva con un motor de búsqueda para explicar el funcionamiento de Attention:

\begin{figure}[H]
\centering
\includegraphics[width=0.75\textwidth]{slides-images/slide-67.jpg}
\caption{Analogía con un motor de búsqueda: el marco Query-Key-Value\protect\footnotemark}
\end{figure}
\footnotetext{Diapositiva 67.}

Supongamos que buscas «deep learning course» (Query) en un motor de búsqueda de videos:
\begin{enumerate}
    \item Cada video en la base de datos tiene una etiqueta o descripción (Key)
    \item El motor de búsqueda calcula la \textbf{similitud} entre la Query y cada Key
    \item Una vez encontrada la Key que mejor coincide, se extrae el contenido de video correspondiente (Value)
\end{enumerate}

\begin{knowledgebox}{La terna Query-Key-Value}
Los tres conceptos centrales del mecanismo de Attention:
\begin{itemize}
    \item \textbf{Query (Q)}: lo que estás buscando, la «pregunta» que necesita responderse en ese momento
    \item \textbf{Key (K)}: la «etiqueta» de cada entrada en la base de datos, el índice que se usa para hacer coincidir
    \item \textbf{Value (V)}: el contenido real asociado a cada Key, la información que se extrae una vez lograda la coincidencia
\end{itemize}
El proceso de Attention consiste en usar la Query para buscar entre las Keys la entrada que mejor coincide, y luego extraer los Values correspondientes.
\end{knowledgebox}

\subsection{Resumen de la sección}

El núcleo del mecanismo de Attention es un proceso de «búsqueda y recuperación»: se calcula la similitud entre la Query y las Keys para determinar a qué partes de la entrada se debe prestar atención, y luego se extrae la información relevante a partir de los Values. Este mecanismo no necesita procesar la secuencia paso a paso y puede modelar directamente la dependencia entre dos posiciones cualesquiera.

\newpage
\section{Self-Attention y Transformer}

\subsection{Los cuatro pasos de Self-Attention}

Self-Attention (autoatención) es una forma especial de Attention: \textbf{Query, Key y Value provienen todos de la misma secuencia de entrada}, es decir, la secuencia calcula atención sobre sus propias posiciones.

\begin{figure}[H]
\centering
\includegraphics[width=0.85\textwidth]{slides-images/slide-70.jpg}
\caption{Los cuatro pasos de Self-Attention\protect\footnotemark}
\end{figure}
\footnotetext{Diapositivas, página 70.}

El proceso completo de cálculo de Self-Attention se divide en cuatro pasos:

\textbf{Paso 1: codificar la información de posición (Positional Encoding)}

Como ya no procesamos la secuencia paso a paso, sino que introducimos todos los elementos a la vez, necesitamos codificar explícitamente la información de posición:

$$\tilde{x}_i = x_i + p_i$$

donde $p_i$ es el vector de codificación posicional, que proporciona un identificador único para cada posición.

\begin{importantbox}{La necesidad de Positional Encoding}
Sin codificación posicional, Self-Attention no podría distinguir entre ``the cat sat on the mat'' y ``the mat sat on the cat'', porque la propia operación de Attention es \textbf{invariante ante permutaciones} (Permutation Invariant) del orden de la entrada. La codificación posicional rompe esta simetría y permite que el modelo perciba las relaciones de orden dentro de la secuencia.
\end{importantbox}

\textbf{Paso 2: extraer Query, Key y Value}

A partir de la entrada ya codificada posicionalmente, se generan las matrices Q, K y V mediante tres capas de transformación lineal distintas:

$$Q = \tilde{X} \cdot W_Q, \quad K = \tilde{X} \cdot W_K, \quad V = \tilde{X} \cdot W_V$$

\begin{itemize}
    \item $W_Q, W_K, W_V$ son tres matrices de pesos que se pueden aprender
    \item la misma entrada $\tilde{X}$, al pasar por transformaciones lineales distintas, captura aspectos diferentes de la información
\end{itemize}

\textbf{Paso 3: calcular los pesos de Attention}

Se usa el \textbf{producto punto escalado} (Scaled Dot-Product) para calcular la similitud entre Query y Key:

$$\text{Attention Score} = \text{softmax}\left(\frac{Q \cdot K^T}{\sqrt{d_k}}\right)$$

\begin{itemize}
    \item $Q \cdot K^T$: el producto punto de Query y Key, que mide la similitud
    \item $\sqrt{d_k}$: el factor de escala ($d_k$ es la dimensión de Key), que evita que valores de producto punto demasiado grandes provoquen el desvanecimiento del gradiente de softmax
    \item $\text{softmax}$: normaliza la similitud a una distribución de probabilidad
\end{itemize}

\begin{figure}[H]
\centering
\includegraphics[width=0.65\textwidth]{slides-images/slide-73.jpg}
\caption{Cálculo de la similitud mediante el producto punto escalado\protect\footnotemark}
\end{figure}
\footnotetext{Diapositivas, página 73.}

\textbf{Paso 4: extraer las características de mayor atención}

Se pondera Value con los pesos de Attention y se suma para obtener la salida final:

$$\text{Output} = \text{softmax}\left(\frac{Q \cdot K^T}{\sqrt{d_k}}\right) \cdot V$$

\begin{importantbox}{La fórmula completa de Self-Attention}
$$\text{Attention}(Q, K, V) = \text{softmax}\left(\frac{Q K^T}{\sqrt{d_k}}\right) V$$

Esta fórmula se llama Scaled Dot-Product Attention y es la unidad de cálculo central de la arquitectura Transformer.
\end{importantbox}

\subsection{Interpretabilidad de los pesos de Attention}

La matriz de pesos de Attention ofrece una forma intuitiva de interpretabilidad: cada elemento $a_{ij}$ de la matriz representa el ``grado de atención'' que la posición $i$ presta a la posición $j$.

\begin{figure}[H]
\centering
\includegraphics[width=0.75\textwidth]{slides-images/slide-75.jpg}
\caption{Visualización de la matriz de pesos de Attention: cuanto más oscuro el color, mayor el grado de atención\protect\footnotemark}
\end{figure}
\footnotetext{Diapositivas, página 75.}

Por ejemplo, en la oración ``He tossed the tennis ball to serve'', la palabra ``ball'' puede tener un peso de atención alto hacia ``tennis'' y ``tossed'', porque estas palabras son las más relevantes para entender el significado de ``ball''.

\subsection{De Attention Head a Transformer}

\begin{figure}[H]
\centering
\includegraphics[width=0.65\textwidth]{slides-images/slide-78.jpg}
\caption{Self-Attention es el bloque de construcción básico de la arquitectura Transformer\protect\footnotemark}
\end{figure}
\footnotetext{Diapositivas, página 78.}

Un cálculo completo de Self-Attention (desde la codificación posicional, pasando por la extracción de Q/K/V, hasta la salida ponderada) constituye un \textbf{Attention Head}.

\subsection{Multi-Head Attention}

Un Transformer real no usa un solo Attention Head, sino \textbf{varios} Attention Head en paralelo: esto es \textbf{Multi-Head Attention}:

$$\text{MultiHead}(Q, K, V) = \text{Concat}(\text{head}_1, \ldots, \text{head}_h) W^O$$

donde cada head usa $W_Q^i, W_K^i, W_V^i$ distintos:

$$\text{head}_i = \text{Attention}(Q W_Q^i, K W_K^i, V W_V^i)$$

\begin{importantbox}{Las ventajas de Multi-Head Attention}
La atención multicabeza permite que el modelo preste atención \textbf{simultáneamente desde distintos ángulos} a diferentes aspectos de la entrada. Por ejemplo, al procesar una oración:
\begin{itemize}
    \item un head puede concentrarse en capturar relaciones de estructura sintáctica
    \item otro head puede concentrarse en la correlación semántica
    \item un tercer head puede atender a la resolución de referencias (quién se refiere a quién)
\end{itemize}
Las salidas de los distintos head se concatenan y se fusionan mediante una transformación lineal, lo que permite al modelo aprovechar de manera integral múltiples patrones de atención.
\end{importantbox}

\subsection{Resumen de la sección}

Self-Attention, mediante cuatro pasos (codificación posicional, extracción de Q/K/V, cálculo de similitud y suma ponderada), logra modelar directamente las dependencias internas de la secuencia. No necesita procesar la secuencia paso a paso, puede calcularse en paralelo y captura directamente dependencias a cualquier distancia. Multi-Head Attention refuerza aún más la capacidad del modelo de entender la entrada desde múltiples ángulos. Estos mecanismos constituyen el núcleo de la arquitectura Transformer.

\newpage
\section{Aplicaciones e impacto del Transformer}

\subsection{Attention Is All You Need}

En 2017, Vaswani et al. publicaron el artículo histórico \textit{Attention Is All You Need}, en el que se propuso la arquitectura Transformer. Esta arquitectura \textbf{prescinde por completo del mecanismo recurrente} y construye el modelo de secuencias basándose únicamente en Self-Attention y capas de retroalimentación directa.

\begin{figure}[H]
\centering
\includegraphics[width=0.6\textwidth]{slides-images/slide-65.jpg}
\caption{«Attention Is All You Need»: el inicio de la era del Transformer\protect\footnotemark}
\end{figure}
\footnotetext{Slides página 65.}

\begin{knowledgebox}{La T de GPT significa Transformer}
Si has usado ChatGPT, entonces la «T» de su nombre representa Transformer. GPT significa Generative Pre-trained Transformer, es decir, un modelo generativo preentrenado basado en la arquitectura Transformer. El Transformer es hoy la arquitectura subyacente de casi todos los grandes modelos de lenguaje.
\end{knowledgebox}

\subsection{Transformer contra RNN: comparación de ventajas}

\begin{table}[H]
\centering
\caption{Comparación entre RNN y Transformer}
\begin{tabular}{lcc}
\toprule
\textbf{Característica} & \textbf{RNN/LSTM} & \textbf{Transformer} \\
\midrule
Modo de procesamiento & Paso a paso en el tiempo (secuencial) & Toda la secuencia en paralelo \\
Dependencias de largo plazo & Difícil (desvanecimiento del gradiente) & Modelado directo (Self-Attention) \\
Eficiencia de entrenamiento & Baja (no se puede paralelizar) & Alta (aprovecha bien la GPU) \\
Información de posición & Implícita (por el orden de los pasos de tiempo) & Explícita (Positional Encoding) \\
Escalabilidad & Limitada & Fuerte (admite escalas muy grandes) \\
\bottomrule
\end{tabular}
\end{table}

\subsection{El amplio alcance de Self-Attention}

El impacto de Self-Attention y del Transformer va mucho más allá del procesamiento de lenguaje natural:

\begin{figure}[H]
\centering
\includegraphics[width=0.85\textwidth]{slides-images/slide-80.jpg}
\caption{Aplicaciones de Self-Attention en distintos campos\protect\footnotemark}
\end{figure}
\footnotetext{Slides página 80.}

\begin{itemize}
    \item \textbf{Procesamiento de lenguaje natural}: grandes modelos de lenguaje como BERT (Devlin et al., NAACL 2019) y GPT (Brown et al., NeurIPS 2020)
    \item \textbf{Secuencias biológicas}: AlphaFold (Jumper et al., Nature 2021) usa Attention para predecir la estructura tridimensional de proteínas; ESMFold (Lin et al., Science 2023)
    \item \textbf{Visión por computadora}: Vision Transformer / ViT (Dosovitskiy et al., ICLR 2020) divide la imagen en una secuencia de parches (Patch) y la procesa con un Transformer, con un desempeño comparable o incluso superior al de las redes neuronales convolucionales
    \item \textbf{Multimodalidad}: los mecanismos de Attention entre modalidades permiten tareas como la generación de imágenes a partir de texto (como DALL-E) y la comprensión de video
\end{itemize}

\begin{importantbox}{La generalidad del Transformer}
La clave del éxito del Transformer está en su diseño arquitectónico \textbf{independiente del dominio} (Domain-agnostic). Sin importar si la entrada es texto, imágenes, secuencias de proteínas o audio, basta con poder representarla como una serie de vectores (Token) para poder procesarla con un Transformer. Esta generalidad ha hecho del Transformer la arquitectura estándar de la era de los «modelos fundacionales» (Foundation Model).
\end{importantbox}

\subsection{Resumen de la sección}

El Transformer transformó por completo el panorama del modelado de secuencias. A través del mecanismo de Self-Attention eliminó la dependencia de estructuras recurrentes, logrando cómputo paralelo eficiente y modelado directo de dependencias de largo alcance. Desde el procesamiento de lenguaje natural hasta la predicción de estructuras de proteínas, y desde la visión por computadora hasta la generación multimodal, el impacto del Transformer está presente en todas partes.

\newpage
\section{Síntesis y ampliación}

\subsection{Hilo conductor central de esta clase}

Esta clase parte del perceptrón más básico y, siguiendo una ruta de evolución clara, construye gradualmente una comprensión completa del modelado de secuencias:

\begin{enumerate}
    \item \textbf{Perceptrón/red feedforward}: puede procesar la entrada de un solo paso de tiempo, pero no puede aprovechar la información histórica
    \item \textbf{RNN}: introduce una ``memoria'' mediante el estado oculto, con lo que logra la transferencia de información entre pasos de tiempo
    \item \textbf{LSTM}: mitiga el problema del desvanecimiento del gradiente mediante un mecanismo de compuertas, mejorando el aprendizaje de dependencias de largo plazo
    \item \textbf{Attention/Transformer}: abandona por completo la estructura recurrente y modela directamente las dependencias globales mediante el mecanismo de autoatención
\end{enumerate}

\begin{importantbox}{El cambio de paradigma de RNN a Transformer}
Esta ruta de evolución refleja un cambio importante de filosofía de diseño en el aprendizaje profundo:
\begin{itemize}
    \item \textbf{Filosofía de RNN}: procesar en orden temporal, manteniendo un estado de memoria comprimido
    \item \textbf{Filosofía de Transformer}: ver todo el contexto de una vez y dejar que la red decida por sí misma en qué enfocarse
\end{itemize}
El éxito de Transformer muestra que, en lugar de que los humanos diseñen la forma en que fluye la información (como la recurrencia), es mejor dejar que la red descubra por sí misma, mediante el aprendizaje, los patrones importantes de los datos.
\end{importantbox}

\subsection{Referencia rápida de conceptos clave}

\begin{table}[H]
\centering
\caption{Tabla de referencia rápida de conceptos clave de esta clase}
\begin{tabular}{ll}
\toprule
\textbf{Concepto} & \textbf{Punto clave} \\
\midrule
Hidden State ($h_t$) & Memoria interna que la RNN transmite entre pasos de tiempo \\
Recurrence Relation & $h_t = f_W(x_t, h_{t-1})$, actualiza de forma recurrente el estado oculto \\
BPTT & Algoritmo que retropropaga el gradiente a lo largo del eje temporal \\
Vanishing Gradient & El gradiente se atenúa con los pasos de tiempo, lo que dificulta el aprendizaje de dependencias de largo plazo \\
LSTM & Variante de RNN que protege el flujo del gradiente mediante un mecanismo de compuertas \\
Embedding & Representación aprendible que mapea símbolos discretos a vectores continuos \\
Positional Encoding & Inyecta información de posición en la entrada del Transformer \\
Self-Attention & Cálculo de atención Query-Key-Value dentro de la secuencia \\
Scaled Dot-Product & $\text{softmax}(QK^T/\sqrt{d_k})V$ \\
Multi-Head Attention & Múltiples grupos de Attention en paralelo, que capturan diversos patrones de atención \\
Transformer & Arquitectura de modelo de secuencias basada por completo en Attention \\
\bottomrule
\end{tabular}
\end{table}

\subsection{Lecturas adicionales}

\begin{itemize}
    \item Vaswani et al., \textit{Attention Is All You Need}, NeurIPS 2017. Artículo pionero de Transformer.
    \item Hochreiter \& Schmidhuber, \textit{Long Short-Term Memory}, Neural Computation 1997. Artículo original de LSTM.
    \item Devlin et al., \textit{BERT: Pre-training of Deep Bidirectional Transformers for Language Understanding}, NAACL 2019.
    \item Brown et al., \textit{Language Models are Few-Shot Learners (GPT-3)}, NeurIPS 2020.
    \item Dosovitskiy et al., \textit{An Image is Worth 16x16 Words: Transformers for Image Recognition at Scale (ViT)}, ICLR 2021.
    \item Sitio oficial del curso MIT 6.S191: \url{http://introtodeeplearning.com}
    \item MIT 6.S191 Lab 1: Deep Learning in Python and Music Generation with RNNs
\end{itemize}

%% --- Fin del contenido --- %%

\end{document}
